\section{Основы эволюционных вычислений и генетических алгоритмов}
\label{sec:ga}

\noindent\hspace{1.25cm}\emph{Раздел используется в ЛР~№5--7.}

\subsection{Эволюционные вычисления}

\emph{Эволюционные вычисления}~--- общее название алгоритмов поиска, оптимизации и обучения,
основанных на формализованных принципах естественного отбора. Эволюционный поиск последовательно
преобразует одно конечное множество промежуточных решений (\emph{популяцию}) в другое: лучшие решения
чаще участвуют в создании новых, новые решения наследуют черты родителей со случайными изменениями, а
худшие постепенно вытесняются. Выделяют четыре основные парадигмы~\cite{skobtsov2008, kureichik2006}:
\begin{itemize}
\item \emph{генетические алгоритмы} (ГА, Дж.~Холланд, 1975~\cite{holland1975})~--- решения кодируются
строками (классически двоичными), основной оператор~--- скрещивание, мутация играет вспомогательную роль;
\item \emph{эволюционные стратегии} (И.~Рехенберг, Х.-П.~Швефель)~--- вещественные векторы, основной
оператор~--- гауссова мутация с самоадаптацией шага;
\item \emph{эволюционное программирование} (Л.~Фогель)~--- эволюция конечных автоматов, затем
вещественных векторов, только мутация;
\item \emph{генетическое программирование} (Дж.~Коза)~--- эволюция программ, представленных деревьями
выражений.
\end{itemize}
Парадигмы различаются представлением решений и набором операторов; современные алгоритмы свободно
комбинируют их идеи. К эволюционным вычислениям примыкают дифференциальная эволюция и
алгоритмы роевого интеллекта (муравьиные, роя частиц).

\emph{Генетический алгоритм}~--- случайно направленный поисковый алгоритм, основанный на механизмах
естественного отбора и генетического наследования. От классических методов оптимизации ГА отличаются
тем, что:
\begin{itemize}
\item работают не с самими параметрами задачи, а с их кодами;
\item ведут поиск из популяции точек, а не из одной точки;
\item используют только значения целевой функции, не требуя производных, непрерывности и выпуклости;
\item применяют вероятностные, а не детерминированные правила перехода.
\end{itemize}
Благодаря этому ГА применимы к многоэкстремальным, негладким и дискретным задачам, но не гарантируют
нахождения глобального оптимума за конечное время и требуют многократных вычислений целевой функции.

\subsection{Основные понятия}

Терминология ГА заимствована из генетики (таблица~\ref{tab:ga-terms}).

\begin{table}[!htb]
\caption{Основные понятия генетических алгоритмов}
\label{tab:ga-terms}
\centering\small
\begin{tabularx}{\textwidth}{>{\raggedright}p{3.6cm}X}
\toprule
Понятие & Значение в ГА \\
\midrule
Особь, хромосома & закодированное допустимое решение задачи, например строка \bits{0110100} \\
Ген & элемент хромосомы: бит, число или номер города \\
Аллель & значение гена (для двоичного кода~--- 0 или 1) \\
Локус & позиция гена в хромосоме \\
Генотип / фенотип & код решения / само решение (значение $x$, маршрут) \\
Популяция & множество из $N$ особей, обрабатываемое на одной итерации \\
Поколение & популяция на очередной итерации $t$ \\
Приспособленность (fitness) & значение функции, которое ГА стремится максимизировать \\
Родители, потомки & особи, участвующие в скрещивании, и особи, полученные в его результате \\
\bottomrule
\end{tabularx}
\end{table}

\emph{Функция приспособленности} (фитнес-функция) $F(x)$ оценивает качество решения. Её выводят из
целевой функции задачи $f(x)$. Если требуется найти максимум положительной функции, можно взять
$F = f$. Для поиска минимума и для функций, принимающих отрицательные значения, применяют
преобразования: $F(x) = C_{\max} - f(x)$, $F(x) = f(x) - f_{\min} + \varepsilon$ (сдвиг на минимум по
текущей популяции, $\varepsilon > 0$ исключает нулевую вероятность отбора) или $F(x) = 1/(1 + f(x))$.
Для методов селекции, использующих только порядок особей (ранжирование, турнир), достаточно сравнивать
значения $f$ со знаком, соответствующим направлению поиска.

\emph{Приспособленность популяции} характеризуют лучшим и средним значениями $F$ по особям; их графики
по поколениям~--- основной инструмент наблюдения за ходом эволюции.

\subsection{Общая структура генетического алгоритма}

Для определения ГА нужно задать: способ кодирования решений, функцию приспособленности, способ создания
начальной популяции, операторы селекции, скрещивания и мутации, стратегию формирования нового
поколения, параметры (размер популяции $N$, вероятности $P_c$ и $P_m$) и критерий останова. Общая схема
показана на рисунке~\ref{fig:ga-flow}.

\begin{figure}[!htb]
\centering
\begin{tikzpicture}[>=Stealth, node distance=5.5mm,
  blk/.style={draw, rounded corners=2pt, minimum width=62mm, minimum height=7.5mm, align=center,
    font=\small, fill=cblue!8},
  dec/.style={draw, diamond, aspect=2.8, inner sep=1pt, align=center, font=\small, fill=cyellow!20},
  term/.style={draw, rounded rectangle, minimum width=30mm, minimum height=7mm, font=\small, fill=clight}]
\node[term] (start) {Начало};
\node[blk, below=of start] (init) {Создание начальной популяции $P(0)$, $t = 0$};
\node[blk, below=of init] (eval) {Вычисление приспособленности особей};
\node[dec, below=of eval] (stop) {Критерий\\останова?};
\node[blk, below=7mm of stop] (sel) {Селекция: выбор родителей};
\node[blk, below=of sel] (cx) {Скрещивание (кроссинговер) с вероятностью $P_c$};
\node[blk, below=of cx] (mut) {Мутация с вероятностью $P_m$};
\node[blk, below=of mut] (eval2) {Вычисление приспособленности потомков};
\node[blk, below=of eval2] (red) {Формирование $P(t+1)$ (сокращение популяции), $t \leftarrow t + 1$};
\node[term, right=18mm of stop] (end) {Лучшая особь};
\draw[->] (start) -- (init);
\draw[->] (init) -- (eval);
\draw[->] (eval) -- (stop);
\draw[->] (stop) -- node[right, font=\small] {нет} (sel);
\draw[->] (stop) -- node[above, font=\small] {да} (end);
\draw[->] (sel) -- (cx);
\draw[->] (cx) -- (mut);
\draw[->] (mut) -- (eval2);
\draw[->] (eval2) -- (red);
\draw[->] (red.west) -- ++(-8mm, 0) |- (stop.west);
\end{tikzpicture}
\caption{Общая схема генетического алгоритма}
\label{fig:ga-flow}
\end{figure}

Критериями останова служат: достижение заданного числа поколений или вычислений функции, отсутствие
улучшения лучшей особи в течение заданного числа поколений, достижение известного оптимального значения
с заданной точностью, вырождение популяции (все особи одинаковы).

\subsection{Простой генетический алгоритм}

\emph{Простой} (классический) ГА, описанный Д.~Голдбергом~\cite{goldberg1989}, использует:
\begin{itemize}
\item двоичное кодирование решений строками фиксированной длины $L$;
\item пропорциональную селекцию (метод <<рулетки>>): особь $i$ выбирается в родительский пул с
вероятностью
\begin{equation}
\label{eq:roulette}
p_i = \frac{F_i}{\sum_{j=1}^{N}F_j};
\end{equation}
\item одноточечный кроссинговер: пары родителей с вероятностью $P_c$ (обычно $0,6$--$0,95$)
обмениваются частями хромосом после случайной точки разреза $k \in \{1, \ldots, L - 1\}$;
\item побитовую мутацию: каждый бит инвертируется с малой вероятностью $P_m$ (обычно $0,001$--$0,01$
или $1/L$);
\item полную замену поколения: потомки заменяют родителей.
\end{itemize}

\emph{Пример.} Найдём максимум $f(x) = x(30 - x)$ на целых числах $x \in [0; 31]$, закодированных пятью
битами; истинный максимум $f(15) = 225$. Пусть начальная популяция из $N = 4$ особей получена случайно
(таблица~\ref{tab:sga-example}). Сумма приспособленностей равна 750, поэтому вероятности отбора
$p_i = f_i / 750$, а ожидаемое число копий $Np_i$.

\begin{table}[!htb]
\caption{Одно поколение простого ГА для $f(x) = x(30 - x)$}
\label{tab:sga-example}
\centering\footnotesize\setlength{\tabcolsep}{4pt}
\begin{tabular}{cccccc|ccc}
\toprule
\multicolumn{6}{c|}{Начальная популяция и селекция} & \multicolumn{3}{c}{После скрещивания} \\
№ & Строка & $x$ & $f(x)$ & $p_i$ & $Np_i$ & Пара, точка & Потомок & $f(x)$ \\
\midrule
1 & \bits{01101} & 13 & 221 & 0,295 & 1,18 & 1--4, $k = 3$ & \bits{011|11} = 15 & 225 \\
2 & \bits{11000} & 24 & 144 & 0,192 & 0,77 & 1--4, $k = 3$ & \bits{100|01} = 17 & 221 \\
3 & \bits{01000} & 8  & 176 & 0,235 & 0,94 & 1--3, $k = 4$ & \bits{0110|0} = 12 & 216 \\
4 & \bits{10011} & 19 & 209 & 0,279 & 1,11 & 1--3, $k = 4$ & \bits{0100|1} = 9 & 189 \\
\midrule
\multicolumn{3}{l}{Сумма / среднее / макс.} & 750 / 187,5 / 221 & & & \multicolumn{2}{r}{} & 851 / 212,8 / 225 \\
\bottomrule
\end{tabular}
\end{table}

Отрезки рулетки: $[0; 0,295)$~--- особь~1, $[0,295; 0,487)$~--- особь~2, $[0,487; 0,721)$~--- особь~3,
$[0,721; 1]$~--- особь~4. Случайные числа $0,12$; $0,27$; $0,55$; $0,90$ выбирают в родительский пул особи
1, 1, 3, 4: лучшая особь получила две копии, худшая (особь~2)~--- ни одной, что соответствует
ожидаемым числам $Np_i$. Пары $(1, 4)$ и $(1, 3)$ скрещиваются в точках $k = 3$ и $k = 4$. Потомок
\bits{01111} ($x = 15$) оказался оптимальным: он объединил начало \bits{011} особи~1 с окончанием
\bits{11} особи~4. Мутация с $P_m = 0,01$ на 20 битах в среднем меняет $0,2$ бита за поколение; в
примере мутаций нет. Средняя приспособленность выросла с $187,5$ до $212,8$ за одно поколение.

\subsection{Представление вещественных решений в двоичной форме}

Для оптимизации функции непрерывной переменной $x \in [a; b]$ отрезок разбивается на $2^L - 1$ равных
частей, и строке $s$ из $L$ бит сопоставляется узел сетки:
\begin{equation}
\label{eq:decode}
x = a + \dec(s)\,\frac{b - a}{2^L - 1}, \qquad \dec(s) = \sum_{j=0}^{L-1} s_j 2^j .
\end{equation}
Шаг сетки $\Delta = (b - a)/(2^L - 1)$ определяет предельную точность: никакой ГА не найдёт решение
точнее $\Delta/2$. Чтобы получить точность $\varepsilon$, длину хромосомы выбирают из условия
\begin{equation}
\label{eq:chrom-length}
2^L - 1 \ge \frac{b - a}{\varepsilon}, \qquad L = \left\lceil \log_2\left(\frac{b - a}{\varepsilon} + 1\right)\right\rceil .
\end{equation}
Например, для $[0; 7]$ и $\varepsilon = 10^{-5}$ получаем $\log_2 700\,001 \approx 19,4$, т.\,е. $L = 20$ и
$\Delta = 6,7\cdot 10^{-6}$. Слишком короткая хромосома делает требуемую точность недостижимой, а каждый лишний
бит вдвое увеличивает пространство поиска без выигрыша в качестве.

При $n$ переменных хромосома~--- конкатенация $n$ строк, каждая из которых декодируется по формуле
(\ref{eq:decode}) на своём отрезке. Следует помнить, что узлы сетки не обязательно содержат точку
оптимума: например, при симметричной области $[-100; 100]$ и чётном числе интервалов нет узла, равного
нулю, так как $(2^L - 1)/2$~--- не целое число. Предел точности ГА тогда определяется значением функции
в ближайшем узле.

\subsection{Код Грея}
\label{sec:gray}

У обычного двоичного кода соседние целые числа могут отличаться многими битами: $7 = \bits{0111}$ и
$8 = \bits{1000}$ различаются всеми четырьмя. Для перехода между ними мутацией нужно одновременно
изменить все биты~--- возникает \emph{<<хэмминговый обрыв>>} (Hamming cliff), мешающий доводке решения.
В \emph{коде Грея} соседние числа отличаются ровно одним битом (таблица~\ref{tab:gray}).

\begin{table}[!htb]
\caption{Двоичный код и код Грея}
\label{tab:gray}
\centering\small
\begin{tabular}{c|cccccccc}
\toprule
Число & 0 & 1 & 2 & 3 & 4 & 5 & 6 & 7 \\
\midrule
Двоичный код & \bits{000} & \bits{001} & \bits{010} & \bits{011} & \bits{100} & \bits{101} & \bits{110} & \bits{111} \\
Код Грея     & \bits{000} & \bits{001} & \bits{011} & \bits{010} & \bits{110} & \bits{111} & \bits{101} & \bits{100} \\
\bottomrule
\end{tabular}
\end{table}

Преобразования (разряды нумеруются от старшего $b_1$ к младшему $b_L$, $\oplus$~--- сложение по модулю~2):
\begin{equation}
\label{eq:gray}
g_1 = b_1, \quad g_i = b_{i-1} \oplus b_i \ (i \ge 2); \qquad
b_1 = g_1, \quad b_i = b_{i-1} \oplus g_i .
\end{equation}
В NumPy прямое преобразование целого числа записывается как \code{g = b \^{} (b \textgreater\textgreater{} 1)}, а обратное для
массива битов~--- как накопленное исключающее ИЛИ \code{np.bitwise\_xor.accumulate(g, axis=1)}, что
равносильно матричной форме $\mathbf{b} = A\mathbf{g} \pmod 2$ с нижней треугольной матрицей $A$ из
единиц. Код Грея сохраняет близость: малое изменение числа требует малого изменения кода, поэтому
мутация чаще порождает близкие решения. Эффект заметен на практике: в примере ЛР~№5 все неудачные
запуски двоичного ГА останавливаются у хэммингового обрыва, а с кодом Грея таких остановок нет.

\subsection{Теория схем и фундаментальная теорема ГА}

Почему ГА работает? Классическое объяснение дано Холландом через понятие \emph{схемы}~---
шаблона в алфавите $\{0, 1, *\}$, где $*$ означает <<любое значение>>. Схема $H = \bits{1**0*}$
описывает все строки длины 5 с единицей в первой позиции и нулём в четвёртой; ей соответствует $2^3 = 8$
строк. Характеристики схемы:
\begin{itemize}
\item \emph{порядок} $o(H)$~--- число фиксированных позиций ($o(\bits{1**0*}) = 2$);
\item \emph{определяющая длина} $\delta(H)$~--- расстояние между первой и последней фиксированными
позициями ($\delta(\bits{1**0*}) = 4 - 1 = 3$).
\end{itemize}
Каждая строка длины $L$ принадлежит $2^L$ схемам, поэтому, оценивая $N$ строк, ГА неявно получает
информацию о гораздо большем числе схем (\emph{неявный параллелизм}; по оценке Холланда
эффективно обрабатывается порядка $N^3$ схем).

Пусть $m(H, t)$~--- число строк популяции, принадлежащих схеме $H$ в поколении $t$, $F(H)$~--- средняя
приспособленность этих строк, $\bar F$~--- средняя приспособленность популяции. Оценим действие операторов
простого ГА.

\emph{Репродукция.} При пропорциональной селекции строка выбирается с вероятностью $F_i / \sum F_j$,
поэтому после $N$ выборов
\begin{equation}
m(H, t + 1) = m(H, t)\,\frac{F(H)}{\bar F}.
\end{equation}
Схемы с приспособленностью выше средней получают больше копий, ниже средней~--- меньше. Если схема
остаётся выше средней на $c\bar F$, её число растёт в геометрической прогрессии: $m(H, t) = m(H, 0)(1 + c)^t$.

\emph{Кроссинговер} разрушает схему, если точка разреза попадает между её фиксированными позициями.
Вероятность этого не превышает $\delta(H)/(L - 1)$, поэтому схема выживает с вероятностью
$p_s \ge 1 - P_c\,\delta(H)/(L - 1)$.

\emph{Мутация} сохраняет схему, если не изменена ни одна из $o(H)$ фиксированных позиций:
$(1 - P_m)^{o(H)} \approx 1 - o(H)P_m$ при $P_m \ll 1$.

Объединяя эти оценки, получаем \emph{фундаментальную теорему ГА} (теорему схем):
\begin{equation}
\label{eq:schema}
m(H, t + 1) \ge m(H, t)\,\frac{F(H)}{\bar F}\left[1 - P_c\,\frac{\delta(H)}{L - 1} - o(H)\,P_m\right].
\end{equation}
\emph{Короткие схемы низкого порядка с приспособленностью выше средней получают экспоненциально
растущее число представителей в следующих поколениях.} Такие схемы называют \emph{строительными
блоками}. \emph{Гипотеза строительных блоков} утверждает, что ГА находит хорошие решения, комбинируя
короткие высокоприспособленные блоки в более длинные~--- именно это произошло в примере
таблицы~\ref{tab:sga-example}, где потомок \bits{01111} собран из блоков \bits{011**} и \bits{***11}.

Теорема схем даёт нижнюю оценку на одно поколение и не учитывает создание новых представителей схемы
кроссинговером, поэтому не является строгим доказательством сходимости. Существуют \emph{обманчивые}
(deceptive) функции, для которых короткие схемы ведут в сторону от оптимума. Тем не менее теорема
объясняет практические правила: гены, совместно влияющие на качество, следует располагать в хромосоме
рядом; вероятность мутации должна быть мала, иначе разрушаются даже хорошие схемы; давление отбора
определяется отношением $F(H)/\bar F$ и зависит от масштаба функции приспособленности.
